% PDF-derived LaTeX transcription. Original manuscript source was unavailable.
% Published PDF SHA-256: 251c871a596ecc98286c82ad63235596a9ab91eb9f2d95ff6a4903cf2f1ad1a0
% Extracted text SHA-256: ea3b26cd95dcbb50b04dc68c345adb4c083a36a00a6f55eec4cc1b8c9a823f40
\documentclass[11pt]{article}
\usepackage[margin=1in]{geometry}
\usepackage{amsmath,amssymb}
\setlength{\emergencystretch}{3em}
\title{Grid-native latent patchification for two-dimensional Maxwell-field reconstruction: a checkpoint-selection-aware retrospective case study}
\author{}
\date{}
\begin{document}
\maketitle
\noindent\textbf{PDF-derived LaTeX transcription.} This editable text follows the PDF extraction. Refer to the published PDF for figures, table layouts, and mathematical notation.

% Extracted chunk 1

Grid-native latent patchification for two-dimensional Maxwell-field
reconstruction: a checkpoint-selection-aware retrospective case
study

\section*{Abstract}

Reconstructing spatial electromagnetic fields from compact learned representations requires
separating representation size, reconstruction error and deterministic constraint enforcement.
We retrospectively examine captured numerical reports for convolutional autoencoders that re-
construct four-frame, ten-channel trajectories on a periodic 96 \ensuremath{\times} 96 grid. A 48-channel latent
on 48 \ensuremath{\times} 48 sites contains 110592 scalars, giving a dynamic-input/latent scalar-count ratio of
3.33. Exact rearrangement into 2 \ensuremath{\times} 2 latent patches produces 576 tokens of width 192 without
further reducing the scalar count. The fixed-vault convolutional model reports raw median/95th-
percentile joint relative- L2 errors of 0.01989/0.06932 on 200 validation trajectories; a separately
trained patchified model with a live transformer reports 0.02941/0.08154 on 256. After a deter-
ministic correction that also uses an external permittivity map, these become 0.01897/0.05641
and 0.02734/0.06700. Three stream-trained configurations also satisfy the recorded corrected
median and 95th-percentile criterion below 0.1, but all five configurations have maximum errors
above 0.1, and correction increases the maximum in two stream cases. Archived pooled-token
baselines have substantially larger reported errors, although unmatched cohorts, optimization
and checkpoint policies preclude causal attribution to pooling alone. The results establish fi-
nite reconstruction feasibility for these grid-native configurations, not losslessness of the learned
bridge or learned electromagnetic conservation. All headline scores are checkpoint-selected val-
idation results rather than independent-test estimates.

\section*{1 Introduction and contributions}

Learned representations of electromagnetic trajectories must preserve spatially distributed infor-
mation while making the representation suitable for downstream computation. Three questions
are easily conflated: whether the learned encoder retains enough information for reconstruction,
whether conversion of a spatial latent to a token sequence introduces additional information loss,
and whether a reconstructed field satisfies a chosen discrete constraint. These questions require
different evidence.
This paper presents a focused retrospective case study of grid-native convolutional latents for
two-dimensional Maxwell-field reconstruction. Its numerical basis is five retained grid-native eval-
uation reports, an archive of pooled-token baseline evaluations, and a separate spectral/resampling
diagnostic. The contribution is an interpretation of those finite results with explicit accounting
for checkpoint selection, scalar counts, sampling and postprocessing. It is not a new transformer
architecture or a claim of globally novel tokenization.
The principal observations are that the retained grid-native models achieve low joint trajectory
reconstruction errors under the recorded validation protocol, and that their spatial latents admit
1

exact patch indexing into a flat sequence. The latter statement concerns only a permutation
of stored entries. Learned encoding, positional arithmetic, transformer processing and decoding
remain distinct operations. Similarly, a small divergence residual after deterministic correction
demonstrates constraint satisfaction by construction, not acquisition of Maxwell dynamics by the
network.

\section*{2 Related work}

Perceiver uses asymmetric cross-attention to map a large input array into a smaller latent array,
followed by latent attention [ 1]. This provides context for the archived pooled-token implementa-
tions examined here. It does not predict a reconstruction floor for electromagnetic fields, and an
implementation-specific failure cannot establish an intrinsic limitation of the architecture.
Vision Transformer processes sequences obtained from flattened image patches, with learned
embeddings and positional information [ 2]. The present case uses patches of a convolutional spatial
latent rather than directly embedding image pixels. Our exactness argument applies to unprojected
patch rearrangement and its inverse; it does not establish invertibility of ViT, positional additions or
a trained transformer. Neither reference supplies evidence for the numerical Maxwell reconstruction
results reported below.

\section*{3 Scientific methods and validation protocol}


\section*{3.1 Task, source fields and metric}

The learned input and output are tensors X \ensuremath{\in} R4\ensuremath{\times}10\ensuremath{\times}96\ensuremath{\times}96 containing
Ex, Ey, Ez, Bx, By, Bz, \ensuremath{\rho}, Jx, Jy, Jz. They are drawn from source arrays with 17 channels; static ma-
terial/source and initial-condition channels are not reconstructed by the learned field latent. The
captured generator describes a periodic, two-dimensional, z-invariant Yee finite-difference time-
domain construction, with four stored frames comprising two consecutive-step pairs. Its implemen-
tation enables JAX float64 computation and stores frames in float32. The arrays retain native Yee
offsets: Ex and Ey occupy (i + 1/2, j) and (i, j + 1/2), Ez and \ensuremath{\rho} occupy nodes, and Bx, By, Bz
occupy (i, j + 1/2), (i + 1/2, j) and (i + 1/2, j + 1/2), respectively. Current components share the
corresponding electric-field offsets. In each implemented step from n to n + 1, the stored frame
contains the updated electric field and charge, the arithmetic average of the old and new magnetic
half-step fields, and the total source-plus-conduction current at n + 1/2. Thus the current stored in
the later frame of a consecutive pair belongs to the interval connecting that pair, not the interval fol-
lowing it. Magnetic divergence uses forward differences, whereas electric-displacement and current
divergence use backward differences on this unit-spacing periodic grid. The joint reconstruction
metric compares these native arrays, not colocated instantaneous continuum fields. These source
definitions describe the intended numerical population, not an independently reproduced solver
validation.
Training estimates a mean \ensuremath{\mu}c and standard deviation \ensuremath{\sigma}c for each dynamic channel from sampled
training data and applies eXc = (Xc \ensuremath{-} \ensuremath{\mu}c)/\ensuremath{\sigma}c. Here, fixed-vault denotes training on stored trajecto-
ries indexed by a fixed manifest, in contrast to live-stream sampling. Fixed-manifest statistics use
up to 256 sampled training trajectories; the stream implementation samples 256 live trajectories
2

Table 1: Retained grid-native configurations. All use a 48 \ensuremath{\times} 48 latent grid and disabled skip paths.
Steps are completed training steps; selected denotes the restored checkpoint step. Counts refer to
the dynamic field representation only.
Model nval C Latent scalars R Steps / selected
Fixed-vault 200 48 110592 3.333 16000 / 15500
Patchified 256 48 110592 3.333 40000 / 39750
Stream A 256 128 294912 1.250 50000 / 48000
Stream B 256 96 221184 1.667 50000 / 49500
Stream C 256 96 221184 1.667 50000 / 35250
for this purpose. Evaluation reverses this normalization before computing, for each trajectory,
r( bX, X ) = \ensuremath{\parallel} bX \ensuremath{-} X\ensuremath{\parallel}2
max(\ensuremath{\parallel}X\ensuremath{\parallel}2, 10\ensuremath{-}12) . (1)
Both norms include all four frames, all ten channels and all grid sites. Thus r is a joint relative- L2
error, not a standard-deviation-normalized error or a mean of channel-wise relative errors. Channel
weighting follows the native numerical scales of the source arrays; the metric is not an SI-calibrated
electromagnetic error. Quantiles and maxima are over trajectories, not grid cells or frames. The
reported trajectory counts do not authenticate statistical independence of their underlying physical
scenes.

\section*{3.2 Grid-native representation and scalar accounting}

The convolutional family uses circular-padding convolutions, residual GroupNorm/SiLU blocks,
strided downsampling, a spatial latent and an upsampling decoder. All five retained reports disable
encoder–decoder skip paths, so the reconstruction cannot bypass the stated spatial bottleneck
through those paths. The fixed-vault report specifies base width 96 and four residual blocks per
configured stage. Its saved training record reports 40 devices and global batch 320; this is not a
matched scaling experiment.
All retained latents have spatial width 48. For C latent channels, the dynamic-input/latent
scalar-count ratio is
R = 4 · 10 · 962
C · 482 . (2)
Consequently, C = 48 gives 110592 latent scalars and R = 10 /3, rather than 2.5. The stream
latents have C = 128 , 96, 96, respectively. Table 1 separates these representation counts from re-
construction scores. These ratios omit model weights, normalization state and the external material
map used by correction. They are neither end-to-end codec bitrates nor like-for-like compression
levels across the configurations.

\section*{3.3 Patch indexing versus learned processing}

For a latent z \ensuremath{\in} RC\ensuremath{\times}g\ensuremath{\times}g and patch size p dividing g, let q = g/p. One indexing convention is
Tiq+j, cp2+ap+b = zc, ip+a, jp+b, (3)
3

where 0 \ensuremath{\leq} i, j < q , 0 \ensuremath{\leq} a, b < p and 0 \ensuremath{\leq} c < C . Each latent entry appears exactly once. Integer
quotient and remainder recover its patch and within-patch coordinates, giving an exact inverse
without numerical arithmetic on the field values. Therefore
N = (g/p)2, D = Cp2, N D = Cg 2. (4)
At C = g = 48 , p = 1 , 2, 4 yields [N, D] = [2304 , 48], [576, 192], [144, 768]. Patch size changes
sequence geometry, not the total number of entries.
The captured bridge performs the broader operation
z\ensuremath{\prime} = P \ensuremath{-}1[U+0000]
F (P (z) + A) \ensuremath{-} A
[U+0001]
, (5)
where P is patch rearrangement, A is a positional array and F is either the identity or a Trans-
formerEncoder. The retained patchified report uses learned positions, p = 2 and a live transformer
at width 192. Subtracting positions after a learned nonidentity map does not invert that map.
Even with identity F , floating-point addition and subtraction need not preserve bits: in float32
round-to-nearest arithmetic, z = 16777216 and A = 1 give (z + A) \ensuremath{-} A = 16777215. The described
Linear(D, 5120) lift is a possible interface, not evidence of executed integration with a trained
ten-billion-parameter model.

\section*{3.4 Deterministic correction}

The magnetic component of the postprocessor is a periodic discrete Fourier Hodge projection. On
an evaluation grid of width g = 96, the unit-spacing forward-difference symbols are
kx = exp(2\ensuremath{\pi}imx/g) \ensuremath{-} 1, k y = exp(2\ensuremath{\pi}imy/g) \ensuremath{-} 1. (6)
For every nonzero Fourier mode, it computes
\ensuremath{\varphi} = kx bBx + ky bBy
|kx|2 + |ky|2 , (7)
bB\ensuremath{\prime}
x = bBx \ensuremath{-} k\ensuremath{\ast}
x\ensuremath{\varphi}, bB\ensuremath{\prime}
y = bBy \ensuremath{-} k\ensuremath{\ast}
y\ensuremath{\varphi}. (8)
Thus kx bB\ensuremath{\prime}
x + ky bB\ensuremath{\prime}
y = 0 in exact arithmetic. The zero mode is unchanged. The implementation uses
float64/complex128 intermediate calculations and casts corrected fields to float32.
The complete correction also changes charge and in-plane electric fields. It first removes each
frame’s spatial charge mean. For the frame pairs (0, 1) and (2, 3), it then replaces the later charge
by the earlier mean-adjusted charge minus \ensuremath{\Delta}t times the backward divergence of the later frame’s
current, using \ensuremath{\Delta}t = 0 .5. It obtains the original relative permittivity map from source channel
11, outside the ten-channel learned representation, and forms edge permittivities by arithmetic
averaging of neighboring node values.
For each frame, the electric correction uses the periodic operator L\ensuremath{\psi} = D\ensuremath{-}
x (\ensuremath{\varepsilon}xD+
x \ensuremath{\psi})+D\ensuremath{-}
y (\ensuremath{\varepsilon}yD+
y \ensuremath{\psi}).
Its right-hand side is the backward divergence of \ensuremath{\varepsilon}E minus the spatially mean-subtracted charge;
the right-hand side’s own mean is also removed to satisfy periodic compatibility. Starting from
zero potential, the conjugate-gradient loop performs 200 iterations, with updates active only when
|\ensuremath{\langle}p, Lp\ensuremath{\rangle}| > 10\ensuremath{-}24 and the squared residual norm exceeds 10\ensuremath{-}16 times its initial value. Inactive up-
dates have zero step and recurrence coefficients; the active recurrence denominator for the squared
residual includes an additive 10\ensuremath{-}30. The solve uses float64 arithmetic and subtracts the forward
4

potential gradient from Ex, Ey before casting to float32. These are implementation specifications,
not demonstrated bounds on Gauss-law or continuity residuals.
Accordingly, we label reconstruction scores “after correction,” rather than attributing them solely
to magnetic projection. The full electric/charge correction has not been independently validated;
the separate algebra checks concern magnetic projection only. The reported magnetic residual is
\ensuremath{\eta}B = RMS(D+
x Bx + D+
y By)q
mean(B2x + B2y ) + 10\ensuremath{-}20
. (9)
This is a discrete unit-spacing residual with its own denominator, distinct from trajectory recon-
struction error.

\section*{3.5 Selection and retrospective verification}

Periodic raw median error on a validation subset selects the checkpoint subsequently restored
for headline evaluation. Final scoring reuses the fixed validation list. The results are therefore
checkpoint-selected validation results, not a fresh independent test. The reports do not authenticate
exact historical row/seed identities or subset nesting across evaluations.
This study is a retrospective reanalysis with no additional physical simulations. No original
checkpoint was loaded, no trained-model inference was repeated and no supplied scientific source
was executed. Separately written NumPy code checked patch indexing and magnetic projection
algebra on synthetic grids of widths 16, 32 and 96. These structural checks support the indexing
and discrete projection interpretation; they do not reproduce original neural predictions, solver
trajectories or full electric/charge correction. Captured sources, full numerical report values and
checking records are retained in Supplementary record S1.

\section*{4 Experimental results}


\section*{4.1 Grid-native reconstruction}

Table 2 gives the canonical restored-checkpoint scores. The fixed-vault model has the lowest re-
ported median and 95th percentile among the five retained grid-native configurations. Its raw
median/95th percentile are 0.01989/0.06932, changing to 0.01897/0.05641 after correction. The
patchified live-transformer model reports 0.02941/0.08154 raw and 0.02734/0.06700 after correc-
tion. This establishes that low reconstruction error was attained in the recorded patchified config-
uration. Because training duration, configuration and evaluation counts differ, it does not measure
an isolated transformer-bridge penalty or establish a zero-cost bridge theorem.
All five corrected medians and 95th percentiles are below the recorded 0.1 criterion. All five
maxima remain above 0.1. Correction lowers p95 in each configuration, but raises the maximum
from 0.15079 to 0.16546 in Stream B and from 0.17717 to 0.18983 in Stream C. A distributional
improvement therefore does not imply improvement for every trajectory. Stream C also illustrates
why raw and corrected endpoints must remain separate: its raw p95 exceeds 0.1, whereas its
corrected p95 does not.
The saved corrected magnetic-divergence residuals cluster near 3.6 \ensuremath{\times} 10\ensuremath{-}8. The independent
synthetic calculation gives float64 relative divergence of approximately 2.3 \ensuremath{\times} 10\ensuremath{-}16 to 3.4 \ensuremath{\times} 10\ensuremath{-}16,
rising to approximately 3.6 \ensuremath{\times} 10\ensuremath{-}8 to 3.8 \ensuremath{\times} 10\ensuremath{-}8 after float32 requantization. This agreement is
5

Table 2: Joint trajectory relative- L2 errors for the selected checkpoints. The complete correction
uses the original permittivity map and modifies charge and electric fields as well as magnetic fields.
Full-precision summaries are retained in Supplementary record S1.
Raw After correction
Model p50 p95 Max p50 p95 Max
Fixed-vault 0.01989 0.06932 0.12598 0.01897 0.05641 0.11353
Patchified 0.02941 0.08154 0.16546 0.02734 0.06700 0.15266
Stream A 0.03079 0.08743 0.14003 0.02913 0.07436 0.13234
Stream B 0.03229 0.09474 0.15079 0.03089 0.08087 0.16546
Stream C 0.03259 0.10027 0.17717 0.03155 0.09303 0.18983
consistent with the projection algebra and storage precision. It is not an independent verification
of the original inference results.

\section*{4.2 Stream sampling}

The three stream configurations record 57582, 46151 and 76982 distinct slot/seed/sample-identifier
tuples from 300006, 200004 and 400008 reads, respectively. Their recorded histories show at most
6000 live slots. These observations support repeated consumption of a recycling slot store, but
identity tuples do not prove distinct tensor content or independent physical experiments.
In the inspected implementation, array files and metadata sidecars are replaced and read sep-
arately. The histories therefore do not establish that each array and its metadata describe the
same historical snapshot. The claimed 6 GB bound on storage or RAM use was not independently
measured. The three configurations are not established independent-seed replications, and differ-
ences in capacity and data consumption prevent attributing their ordering solely to width or data
freshness.

\section*{4.3 Archived baseline outcomes}

The retained baseline archive comprises twenty fleet and seven robust pooled-token evaluation
reports from 37 listed matching directories; ten directories lacked evaluation reports. Complete
historical campaign coverage is therefore not established, and no outcomes are assigned to those
ten directories. Fleet and robust are archive labels, not demonstrated scientific properties. The
robust records describe a training divergence guard that can restore a validation-selected checkpoint
and reduce the learning rate; their headline scores identify restored validation-selected checkpoints.
The fleet reports have no explicit best-checkpoint field. These labels do not define a controlled
architectural or training comparison. Evaluation counts are 96, 128 or 300, unlike the later 200/256
evaluations. None of the retained baseline corrected medians or p95 values meets 0.1.
The best recorded raw median is 0.1623137072 and the best corrected median is 0.1607412472,
both from the same 300-trajectory fleet record. Its raw/corrected p95 values are 0.45944/0.43122. A
stable robust configuration reports raw/corrected medians 0.22339/0.22179 on 96 trajectories. Se-
vere outcomes are also present: one fleet report has raw median/p95/max 37.783/233.024/581.431,
increasing after correction to 58.848/381.729/938.642. Another robust configuration remains near
unit error, with raw/corrected medians 0.99940/0.99924. A divergence-guard event therefore cannot
by itself be interpreted as successful reconstruction.
6

These records establish a finite empirical implementation comparison: the retained later grid-
native configurations attain much smaller reported errors than the archived pooled-token imple-
mentations under their respective protocols. They do not establish a mathematical representation
floor or isolate pooling as the cause. Model capacity, optimization, data cohorts and checkpoint
selection are confounded.

\section*{4.4 Spectral and resampling reference}

The separate diagnostic aggregates 192 two-dimensional time-slice fields, each containing ten chan-
nels, rather than evaluating four-frame trajectories. Its normalized-spectrum energy fraction out-
side the per-axis mode- 16 box is 0.1106126670; the distinct radial-cutoff fraction is 0.1167919636.
This shows appreciable normalized spectral content beyond that box but does not determine the
error of a learned encoder.
Bilinear downsampling to width 48 and upsampling to 96, with aligned-corner interpolation
disabled, yields median/p95 0.0848936476/0.2486859232 after denormalization. The median is
below 0.1, but the tail is not. The width- 96 roundtrip has error near 2 \ensuremath{\times} 10\ensuremath{-}8. These are specific
resampling constructions, not optimality bounds. Their per-slice population and metric granularity
also prevent treating their quantiles as a paired comparison with the trajectory-model scores.

\section*{5 Discussion and limitations}

The scientific significance of this case is the separation of spatial representation learning from
token indexing. Retaining a grid-native latent permits a token interface whose rearrangement is
exactly reversible while leaving the encoder free to learn features before downsampling. A width- 48
latent is therefore not equivalent to bilinear restriction of each physical channel to a width- 48 field.
The observed reconstruction results are consistent with useful learned features, but unmatched
diagnostic and model populations do not quantify that mechanism causally.
Scalar accounting further changes the interpretation. The fixed and patchified cases reduce the
dynamic scalar count by 3.33, while the wider stream latents reduce it by only 1.25 or 1.67. A
lower error at a different latent size is not a matched compression result. Moreover, correction
receives side information that a complete field codec would have to supply separately. Neither
patch indexing nor the small corrected divergence residual makes the full learned system lossless
or removes its computational and physical-validation costs.
Constraint satisfaction and reconstruction fidelity are related but different endpoints. The mag-
netic Hodge operation removes a discrete longitudinal component by construction. The evaluated
complete correction changes additional quantities, and the observed increases in maximum error
show that its aggregate benefit is not uniform. The retained magnetic residuals do not establish
learned conservation, full Gauss-law or continuity correctness, faithful time stepping, continuum
convergence or recovered electromagnetic dynamics.
The main statistical limitation is reuse of validation data for checkpoint selection and scoring.
No fresh test cohort or authenticated independent-seed replication supports uncertainty estimates.
We therefore do not infer confidence intervals from trajectories, frames, grid cells or the number of
archived reports. Exact training-stream/validation disjointness and physical distinctness are not
established, so memorization cannot be ruled out. Scalar report summaries also cannot reproduce
7

the original quantiles without per-trajectory errors or prediction arrays. Current source captures
explain the reported methods but are not exact historical environment attestations.
These limitations bound the result to a finite reconstruction case study. A stronger compar-
ison would require matched latent budgets, authenticated common evaluation trajectories, inde-
pendently selected test data, controlled bridge ablations and reconstruction of the full correction
diagnostics.

\section*{6 Conclusion}

The captured primary reports establish low checkpoint-selected validation errors for five grid-native
Maxwell-field reconstruction configurations. A 48 \ensuremath{\times} 48 \ensuremath{\times} 48 spatial latent can be rearranged exactly
into 576 width-192 tokens, and one retained model processes those tokens with a live transformer
while maintaining a corrected median/p95 of 0.02734/0.06700. This is finite evidence for the tested
representation and interface, not proof of a lossless learned bridge or integration with a larger
trained model.
Interpretation must account for the representation sizes and failure cases: dynamic scalar-count
ratios vary from 1.25 to 3.33, correction consumes an external permittivity map, every retained
model has maximum error above 0.1, and two stream maxima worsen after correction. Archived
baseline failures motivate grid-native representations empirically, without demonstrating a univer-
sal pooling limitation. Exact indexing, selected validation accuracy and constructed constraint
satisfaction should remain separately reported.

\section*{7 Data and Code A vailability}

Captured training, patchification, generator, streaming, correction and spectral-diagnostic source
files, together with numerical evaluation reports, exist in a private local export. Supplementary
record S1 contains the captured records, full-precision summaries and independent checking in-
puts/results. Access established for this study is limited to that private export; no public archive,
external distribution grant or reviewer-access mechanism is established.
The inspected source reads trajectory arrays from private manifest and reservoir locations. Orig-
inal target/prediction arrays and checkpoint contents were not collected or loaded for this re-
analysis; quantiles were not reproduced and inference was not replayed. Independent checking
scripts and their saved outputs exist for report arithmetic, static-definition consistency and syn-
thetic patch/projection algebra. These artifacts support the bounded checks described here, not
full reproduction of the original training or solver campaign.

\section*{References}

[1] A. Jaegle, F. Gimeno, A. Brock, A. Zisserman, O. Vinyals, and J. Carreira. Perceiver: Gen-
eral Perception with Iterative Attention. Proceedings of the 38th International Conference on
Machine Learning, PMLR 139, 4651–4664 (2021). https://proceedings.mlr.press/v139/
jaegle21a.html.
8

[2] A. Dosovitskiy et al. An Image is Worth 16x16 Words: Transformers for Image Recognition
at Scale. International Conference on Learning Representations (2021). arXiv:2010.11929v2.
https://arxiv.org/abs/2010.11929v2.
9

\end{document}
